% sec_conclusiones.tex
% Sección 5: Conclusiones
\section{Conclusiones}

\begin{enumerate}
    \item \texttt{FastBranch} alcanza una latencia de 4.57 ns/op, convergiendo exitosamente al límite teórico físico (3.90 ns/op) con un \textit{speedup} empírico medio de $3.02\times$ (IC 95\%: $[2.90\times, 3.14\times]$, Cohen's $d = 1.82$, $p_{\text{boot}} < 0.0001$).
    \item La dispersión y \textit{jitter} temporal se reduce significativamente ($\sigma \approx 0.87$ ns frente a 1.88 ns en \texttt{switch}), dotando al motor de \textit{branching} de un determinismo crítico para la microestructura de trading.
    \item El uso de punteros crudos acoplados con aislamiento físico de caché (\texttt{alignas(64)}) y semántica \textit{lock-free} en C++17 elimina la penalización por \textit{false sharing} inter-hilo, permitiendo escalabilidad multi-core sin contención en el \textit{hot path}.
    \item Se valida que la técnica semi-estática es contraproducente frente al predictor de hardware BPU únicamente cuando la condición tiene una estabilidad superior al 99.9\% ($f_{\text{set}} \le 1/10^4$), confirmando su superioridad en escenarios de mercado volátiles e impredecibles.
\end{enumerate}
